\documentclass[11pt]{article}
\usepackage[margin=1in]{geometry}
\usepackage{amsmath,amssymb,amsthm}
\usepackage[hidelinks]{hyperref}
\newtheorem{theorem}{Theorem}
\newtheorem{corollary}{Corollary}
\title{Parity of fixed points in triangular-divisor sequences\\
\large A solution to MathDB 343324}
\author{Principia Math In Progress}
\date{18 September 2026}
\begin{document}
\maketitle
\begin{abstract}
We prove the parity conjecture for the sequences of Navarrete and
Orellana. The obstruction is already present in the defining
divisibility condition: at an even index, the index does not divide
the prescribed odd multiple of a triangular number. Primality and
the greedy selection rule are unnecessary for this obstruction.
For comparison, every even positive multiplier gives the identity
sequence under the greedy rule.
\end{abstract}

\section*{Task 1: the original problem and its full solution}
For an odd prime $p$, put
\[
 q_p(n)=\sum_{j=0}^{n-1}jp=\frac{pn(n-1)}2.
\]
The source \cite{NO} defines $a_p(1)=1$ and, for $n\geq2$,
chooses $a_p(n)$ to be the least positive divisor of $q_p(n)$
not among $a_p(1),\ldots,a_p(n-1)$. Its Conjecture 3.1,
on page 4, asserts that $a_p(n)=n$ can occur only at odd $n$.

\begin{theorem}\label{main}
Let $m$ be an odd positive integer. If $n$ is a positive integer
and $n$ divides $mn(n-1)/2$, then $n$ is odd.
Consequently, every fixed point of $A(p)$ is odd for every
odd prime $p$.
\end{theorem}
\begin{proof}
The divisibility assumption says that, for some integer $b$,
\[
 \frac{mn(n-1)}2=nb.
\]
Since $n>0$, cancellation gives $m(n-1)=2b$.
The integer $m$ is odd, so $n-1$ is even and $n$ is odd.
For the application, the defining rule always gives
$a_p(n)\mid q_p(n)$. Substituting $a_p(n)=n$ supplies exactly
the preceding divisibility assumption. This also includes $n=1$,
since every nonzero integer divides zero.
\end{proof}

An equivalent explicit remainder is useful for checking the
obstruction. Writing $m=2\ell+1$ and $n=2k$, where $\ell\geq0$
and $k\geq1$, gives
\begin{equation}\label{remainder}
 \frac{m n(n-1)}2
 =2k\bigl((2\ell+1)k-\ell-1\bigr)+k
 \equiv k\pmod{2k}.
\end{equation}
The remainder is strictly between zero and the modulus. No
information about earlier sequence entries is required.

\begin{corollary}\label{dichotomy}
For each integer $m\geq2$, define the greedy triangular-divisor
sequence with multiplier $m$ by the same rule. It is well defined.
If $m$ is odd, all its fixed points are odd. If $m$ is even,
then $a_m(n)=n$ for every $n\geq1$.
\end{corollary}
\begin{proof}
For $n\geq2$, the positive numbers $q_m(n)$ strictly increase
and $q_m(2)=m>1$. Every previous selected divisor is at most
its corresponding $q_m(j)$, apart from the initial value $1$.
Thus $q_m(n)$ itself has not been selected and is an available
divisor. This proves existence at each step. The odd case follows
from Theorem~\ref{main}. If $m$ is even, then
$q_m(n)=n\bigl((m/2)(n-1)\bigr)$, so $n$ is an available divisor
whenever the preceding values are $1,\ldots,n-1$.
Every smaller positive integer has already been selected;
therefore the greedy rule selects $n$. Induction starts at $a_m(1)=1$.
\end{proof}

\paragraph{Scope and attribution.}
This resolves the original all-prime parity statement, without
changing its indexing. The result does not say that odd indices
are fixed points or that fixed points are prime. The source's
Table 4 already records the identity sequence for multipliers
$2,4,6$; Corollary~\ref{dichotomy} includes these observations.
The later paper \cite{AA} proves a separate result on prime fixed
points below $2p$, along with other assertions. Its results are
credited, and are not used in the proof above.

\section*{Task 2: dependence on odd composite multipliers}
For odd composite $m$, determine when the complete greedy sequence
$A(m)$ agrees with $A(p)$ for a prime divisor $p$ of $m$.
The source reports agreement for $m=9$ and $p=3$.
Establish a criterion or locate the first disagreement for other
prime powers and products of distinct primes.

\section*{Task 3: distribution of composite fixed points}
For a fixed odd prime $p$, study
\[
 C_p(X)=\#\{n\leq X:n\text{ is composite and }a_p(n)=n\}.
\]
Determine whether this count has a natural density and obtain
bounds that use the exclusion of even indices. The existence of
composite fixed points is already documented in the sources;
their distribution is the proposed question.

\section*{Task 4: prime fixed points beyond the known range}
Theorem 5.1 of \cite{AA} gives $a_p(\ell)=\ell$ for odd primes
$\ell\ne p$ with $\ell<2p$. Determine which additional primes
outside this range are fixed points. Seek a criterion involving
the prior availability of divisors of $p\ell(\ell-1)/2$ and
quantitative bounds on the exceptional primes.

\section*{Task 5: other integer-valued input polynomials}
Replace $mn(n-1)/2$ by another increasing integer-valued polynomial
$Q(n)$ and choose the least unused positive divisor at each step.
Classify the residue classes forbidden to fixed points by the
necessary condition $n\mid Q(n)$. Then determine when these local
obstructions account for all missing fixed-point classes.

\paragraph{Verification and status.}
The proof is elementary and does not depend on computation.
The accompanying script checks the exact remainder identity,
reproduces an independently recorded sequence, and checks finite
prefixes for odd and even multipliers. The literature audit
documents the current primary, the related 2022 paper, a bounded
prior-claim search, and live MathDB checks. It is not an exhaustive
novelty guarantee. Review was by the authoring agent only.

\begin{thebibliography}{9}
\bibitem{NO} E. Navarrete and D. Orellana,
\emph{Finding Prime Numbers as Fixed Points of Sequences},
arXiv:1907.10023v1 (2019), 9 pages.
\href{https://arxiv.org/abs/1907.10023}{arXiv record}.
\bibitem{AA} J. A. Julca Avila and G. Silva de Andrade,
\emph{About the theorem that partially solves the
Navarrete--Orellana Conjecture}, Notes on Number Theory and
Discrete Mathematics \textbf{28}(2) (2022), 302--317.
\href{https://doi.org/10.7546/nntdm.2022.28.2.302-317}{doi:10.7546/nntdm.2022.28.2.302-317}.
\end{thebibliography}
\end{document}
